The method is defined for one product class at a time. Fix a class of
mass-market power-electronic unit installed on board (cabin lighting drivers,
cabin or public-area supplies, charging outlets, fan-coil drives) and state
the decision as follows.

\subsection{Specification of the Decision}
\subsubsection{Installed Population}
$N$ units of a known product in one zone or deck of
one ship, from an inventory giving product, manufacturer, installation date and
location. The inventory must be complete, since a miscounted population biases the
exposure, and hence the rate, in direct proportion.

\subsubsection{Failure Event}
A functional failure of the unit leading to its
replacement, as recorded by the maintenance system.

\subsubsection{Exposure}
Powered unit-hours of that population during the mission,
obtained from the installed count and the duty schedule of the zone.

\subsubsection{Mission Horizon}
The interval between the calls at which spares can be
replenished, typically days to weeks. The decision is taken before sailing and
the horizon is bounded and known.

\subsubsection{Data Holders}
The ship holds its own maintenance and environmental
records; the operator holds those of its ships; an integrator or manufacturer
may hold qualification data. None is assumed to release unit-level records.

\subsubsection{Decision}
The vector of stocks to embark. A zone carries $P$ product
references, and a spare of one reference does not replace a failure of another,
so the decision is $\mathbf s=(s_1,\dots,s_P)$ and not a single number. Two
coverage criteria must be distinguished, because they are different problems:
\begin{align}
\text{joint:}\quad &\Pr\big(Y_p\le s_p\ \text{for all } p\big)\ \ge\ \gamma,
\label{eq:joint}\\
\text{aggregate:}\quad &\Pr\Big(\textstyle\sum_p Y_p\le S\Big)\ \ge\ \gamma .
\label{eq:agg}
\end{align}
Criterion~\eqref{eq:joint} is the one the shipboard decision requires: the
mission is covered only if every reference is covered.
Criterion~\eqref{eq:agg} is legitimate only where the references are genuinely
interchangeable, as for a common module or a kit drawn against any of them, and
it is the criterion under which risk pools across references.
Section~\ref{sec:multi} treats both and shows that they respond to the
dependence structure in opposite directions. Where volume or cost limits the
stock, \eqref{eq:joint} becomes a constrained problem,
$\min\sum_p c_ps_p$ subject to the coverage constraint, which the predictive
distribution supplies but which is not pursued further here.

\subsubsection{Covariates Available at the Decision}
Product and manufacturer, the age
of the installed population, its zone, the planned mission length, and
summaries of the previous mission's operating conditions (zone temperature and
the supply quality of the ship's microgrid). Anything measured
during the mission being predicted is not available and must not enter the
predictor. This is a restriction on the model, not a description of what a fleet
happens to record, and it is stated here because a predictor validated with
contemporaneous covariates does not validate this decision.

\subsection{Mechanisms to Be Declared}
\label{sec:mechanisms}
The count model of Section~\ref{sec:model} is a statement about a well-defined
event in a well-defined population over a bounded horizon. Four mechanisms
determine whether shipboard records satisfy that definition, and a fleet must
declare each before its data are modelled. They are the usual reason maintenance
records cannot be used as they stand, and their consequences differ.

\subsubsection{Right Censoring at the Horizon}
A mission ends at a fixed time, so a unit
that has not failed is censored rather than surviving indefinitely. In a count
model over a bounded horizon the exposure offset handles this, the population
contributing the unit-hours it actually accumulated. Populations recording no
failure must remain in the fit; excluding them selects on the outcome and biases
the holder term downwards.

\subsubsection{Preventive Replacement}
A unit replaced on schedule, or opportunistically while a compartment is open, has
not failed. Such replacements must be flagged separately, excluded from the event
count, and their unit-hours deducted from the exposure: recorded as failures they
inflate the rate, and with exposure undeducted they deflate it. A fleet unable to
separate the two categories cannot estimate this model, and the extent of the
ambiguity should be reported as a sensitivity.

\subsubsection{Units Removed for Other Reasons}
Refits, refurbishment, obsolescence and cannibalisation remove units mid-horizon,
contributing the exposure accrued before removal. Removals correlated with
condition, such as a unit withdrawn on suspicion rather than after failure, are
informative and require a declared rule and a sensitivity analysis under both
classifications.

\subsubsection{Dependence Between Failures}
The count model treats failures within a
population as conditionally independent given the latent rate, with all shared
variation carried by the random effects. On a ship that assumption can fail in
ways that matter. A common cause, such as a supply transient, a cooling loss, a
flooding or vibration event or a single bad production lot installed together,
produces simultaneous failures that no smooth latent rate reproduces and inflates
the upper tail beyond a mixed Poisson model; cascading failures, where the loss of
one unit stresses its neighbours, do the same. Such events must be identified in
the records and either modelled separately or excluded with the exposure they
affect, and their frequency reported. This is the assumption most likely to be
violated on board and the one a first campaign should test.

\subsection{Divided Evidence}
Each ship sees few failures of a given product in one mission; a fleet sees many
more, but its records sit with different parties, and commercial sensitivity,
data-protection undertakings and contractual limits mean unit-level records are
commonly not pooled. The consequence is not merely a smaller sample: a predictor
fitted without a holder term is over-confident for a holder it has not seen, while
one fitted on a single ship has almost no information.
Section~\ref{sec:model} addresses the first problem and Section~\ref{sec:fed} the
constraint that produced it.
